% GENERATED FILE. Edit canonical lessons, then run npm run generate:books.
% canonical-source-hash: fnv1a64:ecf3ca9d3379c077
% canonical-lessons: RU-C443-brodit, RU-C443-progulivatsya, RU-C443-kvartal, RU-C443-prigorod, RU-C443-mostovaya

\chapter{A stroll round town}
\label{ch:ru-a-stroll-round-town}
\hlchaptermodality{443}

\begin{chapteropening}
\textbf{By the end of this chapter:} \emph{I can talk about strolling round a town: wandering, taking a stroll, a city block, a suburb and the roadway.}

Complete the last lesson of chapter 443: I can talk about strolling round a town: wandering, taking a stroll, a city block, a suburb and the roadway.
\end{chapteropening}

\section[\texorpdfstring{\ru{бродить} (brodít)}{бродить (brodít)}]{\ru{бродить} (brodít) — to wander, to stroll}

\label{lesson:RU-C443-brodit}



\begin{quote}
Before the new one: say the Russian for a sight, then the Russian for to look round.
\end{quote}

\subsection*{You'll want to know: \ru{бродить}}
\textbf{\ru{бродить}} (\emph{brodít}) — \textquotedblleft{}to wander, to stroll\textquotedblright{}.

\begin{quote}
\ru{Мы} \ru{любим} \ru{бродить} \ru{по} \ru{старым} \ru{улицам}.
\end{quote}

\begin{quote}
\emph{My lyúbim brodít po stárym úlitsam.}
\end{quote}

\begin{quote}
We love wandering the old streets.
\end{quote}

Imperfective. Present \emph{\ru{брожу}, \ru{бродишь}}. Of wine: to ferment.

\begin{cousinweb}
\emph{\ru{бродить}} is inherited through Old East Slavic \emph{broditi} from Proto-Slavic \emph{broditi} (EtymDB).
\end{cousinweb}

\begin{grammarlens}[title={What you've built}]
One more word for telling, explaining and describing at B1.
\end{grammarlens}

\subsection*{Your turn}
\begin{itemize}
\raggedright
  \item \emph{Say it:} \emph{brodít}
  \item \emph{Say it:} \emph{brodít}, once more
  \item \emph{Say it:} a sentence of your own with \emph{brodít}
  \item \emph{Recall:} say the Russian for a sight, then the Russian for to look round, then say \emph{brodít} again
\end{itemize}

\begin{tcolorbox}[breakable,colback=teal!4,colframe=teal!35!black,title={Before you move on}]
What does \ru{бродить} mean? (\textquotedblleft{}To wander, to stroll\textquotedblright{}.) Say it once more.
\end{tcolorbox}
\section[\texorpdfstring{\ru{прогуливаться} (progúlivatsya)}{прогуливаться (progúlivatsya)}]{\ru{прогуливаться} (progúlivatsya) — to take a stroll, to stroll}

\label{lesson:RU-C443-progulivatsya}



\begin{quote}
Before the new one: say the Russian for to look round, then the Russian for to wander.
\end{quote}

\subsection*{You'll want to know: \ru{прогуливаться}}
\textbf{\ru{прогуливаться}} (\emph{progúlivatsya}) — \textquotedblleft{}to take a stroll, to stroll\textquotedblright{}.

\begin{quote}
\ru{Вечером} \ru{люди} \ru{прогуливаются} \ru{в} \ru{парке}.
\end{quote}

\begin{quote}
\emph{Vécherom lyúdi progúlivayutsya v párke.}
\end{quote}

\begin{quote}
In the evening people stroll in the park.
\end{quote}

Imperfective; perfective \emph{\ru{прогуляться}}.

\begin{cousinweb}
\emph{\ru{прогуливаться}} is built from \emph{\ru{про}-}, \emph{\ru{гулять}} 'to go for a walk' and \emph{-\ru{ся}}. EtymDB does not list the verb, and lists \emph{\ru{гулять}} without a source.
\end{cousinweb}

\begin{grammarlens}[title={What you've built}]
One more word for telling, explaining and describing at B1.
\end{grammarlens}

\subsection*{Your turn}
\begin{itemize}
\raggedright
  \item \emph{Say it:} \emph{progúlivatsya}
  \item \emph{Say it:} \emph{progúlivatsya}, once more
  \item \emph{Say it:} a sentence of your own with \emph{progúlivatsya}
  \item \emph{Recall:} say the Russian for to look round, then the Russian for to wander, then say \emph{progúlivatsya} again
\end{itemize}

\begin{tcolorbox}[breakable,colback=teal!4,colframe=teal!35!black,title={Before you move on}]
What does \ru{прогуливаться} mean? (\textquotedblleft{}To take a stroll, to stroll\textquotedblright{}.) Say it once more.
\end{tcolorbox}
\section[\texorpdfstring{\ru{квартал} (kvartál)}{квартал (kvartál)}]{\ru{квартал} (kvartál) — a block (in a city); a quarter (of a year)}

\label{lesson:RU-C443-kvartal}



\begin{quote}
Before the new one: say the Russian for to wander, then the Russian for to take a stroll.
\end{quote}

\subsection*{You'll want to know: \ru{квартал}}
\textbf{\ru{квартал}} (\emph{kvartál}) — \textquotedblleft{}a block (in a city); a quarter (of a year)\textquotedblright{}.

\begin{quote}
\ru{Аптека} \ru{в} \ru{соседнем} \ru{квартале}.
\end{quote}

\begin{quote}
\emph{Aptéka v sosédnem kvartále.}
\end{quote}

\begin{quote}
The chemist's is in the next block.
\end{quote}

Masculine. \emph{\ru{Первый} \ru{квартал}} — the first quarter of the year.

\begin{cousinweb}
\emph{\ru{квартал}} was borrowed from German \emph{Quartal} (EtymDB), a word built on Latin \emph{quartus} 'fourth' (Lewis \& Short).
\end{cousinweb}

\begin{grammarlens}[title={What you've built}]
One more word for telling, explaining and describing at B1.
\end{grammarlens}

\subsection*{Your turn}
\begin{itemize}
\raggedright
  \item \emph{Say it:} \emph{kvartál}
  \item \emph{Say it:} \emph{kvartál}, once more
  \item \emph{Say it:} a sentence of your own with \emph{kvartál}
  \item \emph{Recall:} say the Russian for to wander, then the Russian for to take a stroll, then say \emph{kvartál} again
\end{itemize}

\begin{tcolorbox}[breakable,colback=teal!4,colframe=teal!35!black,title={Before you move on}]
What does \ru{квартал} mean? (\textquotedblleft{}A block (in a city); a quarter (of a year)\textquotedblright{}.) Say it once more.
\end{tcolorbox}
\section[\texorpdfstring{\ru{пригород} (prígorod)}{пригород (prígorod)}]{\ru{пригород} (prígorod) — a suburb}

\label{lesson:RU-C443-prigorod}



\begin{quote}
Before the new one: say the Russian for to take a stroll, then the Russian for a block.
\end{quote}

\subsection*{You'll want to know: \ru{пригород}}
\textbf{\ru{пригород}} (\emph{prígorod}) — \textquotedblleft{}a suburb\textquotedblright{}.

\begin{quote}
\ru{Мы} \ru{живём} \ru{в} \ru{пригороде} \ru{и} \ru{ездим} \ru{в} \ru{город} \ru{на} \ru{поезде}.
\end{quote}

\begin{quote}
\emph{My zhivyóm v prígorode i yézdim v górod na póyezde.}
\end{quote}

\begin{quote}
We live in a suburb and go into town by train.
\end{quote}

Masculine; the stress is on \emph{\ru{при}-} (\emph{prí-}).

\begin{cousinweb}
\emph{\ru{пригород}} is built from \emph{\ru{при}-} 'near' and \emph{\ru{город}} 'a town', which is inherited from Proto-Slavic \emph{gordŭ} 'an enclosure' (EtymDB).
\end{cousinweb}

\begin{grammarlens}[title={What you've built}]
One more word for telling, explaining and describing at B1.
\end{grammarlens}

\subsection*{Your turn}
\begin{itemize}
\raggedright
  \item \emph{Say it:} \emph{prígorod}
  \item \emph{Say it:} \emph{prígorod}, once more
  \item \emph{Say it:} a sentence of your own with \emph{prígorod}
  \item \emph{Recall:} say the Russian for to take a stroll, then the Russian for a block, then say \emph{prígorod} again
\end{itemize}

\begin{tcolorbox}[breakable,colback=teal!4,colframe=teal!35!black,title={Before you move on}]
What does \ru{пригород} mean? (\textquotedblleft{}A suburb\textquotedblright{}.) Say it once more.
\end{tcolorbox}
\section[\texorpdfstring{\ru{мостовая} (mostováya)}{мостовая (mostováya)}]{\ru{мостовая} (mostováya) — a roadway; a paved street}

\label{lesson:RU-C443-mostovaya}



\begin{quote}
Before the new one: say the Russian for a block, then the Russian for a suburb.
\end{quote}

\subsection*{You'll want to know: \ru{мостовая}}
\textbf{\ru{мостовая}} (\emph{mostováya}) — \textquotedblleft{}a roadway; a paved street\textquotedblright{}.

\begin{quote}
\ru{Машины} \ru{едут} \ru{по} \ru{мостовой}, \ru{а} \ru{люди} \ru{идут} \ru{по} \ru{тротуару}.
\end{quote}

\begin{quote}
\emph{Mashíny yédut po mostovóy, a lyúdi idút po trotuáru.}
\end{quote}

\begin{quote}
The cars drive on the roadway, and people walk on the pavement.
\end{quote}

An adjective used as a noun: \emph{\ru{по} \ru{мостовой}}.

\begin{cousinweb}
\emph{\ru{мостовая}} is built on \emph{\ru{мостить}} 'to pave', from \emph{\ru{мост}} 'a bridge', which is inherited from Proto-Slavic \emph{mostŭ} (EtymDB).
\end{cousinweb}

\begin{grammarlens}[title={What you've built}]
One more word for telling, explaining and describing at B1.
\end{grammarlens}

\subsection*{Your turn}
\begin{itemize}
\raggedright
  \item \emph{Say it:} \emph{mostováya}
  \item \emph{Say it:} \emph{mostováya}, once more
  \item \emph{Say it:} a sentence of your own with \emph{mostováya}
  \item \emph{Recall:} say the Russian for a block, then the Russian for a suburb, then say \emph{mostováya} again
\end{itemize}

\begin{tcolorbox}[breakable,colback=teal!4,colframe=teal!35!black,title={Before you move on}]
What does \ru{мостовая} mean? (\textquotedblleft{}A roadway; a paved street\textquotedblright{}.) Say it once more.
\end{tcolorbox}
